\chapter{Illustrative Analytical and Numerical Results}
\label{ch:model-results}

\section{Purpose and limitations of the calculations}

This chapter demonstrates the behavior of the coupled framework using normalized, reproducible calculations. It does not present experimental observations and does not claim to reproduce the full differential Bethe--Heitler or electron--electron kernel. The calculations are intended to answer three design questions before a high-cost simulation or experiment is performed:
\begin{enumerate}
\item How strongly can high-energy redistribution change radiation-weighted moments?
\item Why does the reduction grow when the emission kernel weights energy more strongly?
\item Why can dose reduction differ from integrated source-power reduction after spectral weighting?
\end{enumerate}
The complete script is included in Appendix~\ref{app:computation} and in the project directory.

\section{Reference distribution}

The reference population is a Maxwell--J\"uttner energy distribution at $\theta=0.40$, corresponding to $k_BT_e\approx\SI{204}{\keV}$. The normalized energy probability density is
\begin{equation}
p(E)\propto\gamma^2\beta e^{-\gamma/\theta},
\qquad \gamma=1+\frac{E}{m_ec^2}.
\end{equation}
A cutoff-and-redistribution model moves all electrons above $E_c$ into a narrow lower-energy kernel centered at $E_r=0.30m_ec^2\approx\SI{153}{\keV}$. The example in \Cref{fig:distribution-cutoff} uses $E_c=1.20m_ec^2\approx\SI{613}{\keV}=3k_BT_e$.

\begin{figure}[H]
\centering
\includegraphics[width=0.78\textwidth]{distribution_cutoff.pdf}
\caption{Normalized Maxwell--J\"uttner reference distribution and a number-conserving cutoff-plus-redistribution example. The deposited low-energy feature is a numerical representation of particles transferred from the superthermal tail.}
\label{fig:distribution-cutoff}
\end{figure}

The model preserves particle number to numerical precision. It does not preserve electron energy because the control mechanism is presumed to transfer the difference to another reservoir. A realistic simulation would track that reservoir explicitly.

\section{Tail fraction and radiation-weighted moments}

The tail number fraction is
\begin{equation}
F_N(E_c)=\int_{E_c}^{\infty}p(E)\,\dd E.
\end{equation}
For the selected relativistic temperature, a cutoff at only a few $k_BT_e$ can contain a non-negligible particle fraction because the relativistic density-of-states factor shifts the energy distribution. This illustrates why a proposed cutoff must be defined in physical energy and not interpreted from nonrelativistic intuition alone.

\begin{table}[H]
\centering
\caption{Illustrative number fractions and radiation-weighted moment ratios at $\theta=0.40$. $M_q$ is defined by Equation~\eqref{eq:moment}. Values are generated by the included script and are not full bremsstrahlung power ratios.}
\label{tab:moment-results}
\begin{tabular}{ccccc}
\toprule
$E_c/k_BT_e$ & Tail number fraction & $M_1/M_{1,0}$ & $M_2/M_{2,0}$ & $M_3/M_{3,0}$ \\
\midrule
1.5 & 0.546 & 0.322 & 0.0816 & 0.0181 \\
2.0 & 0.409 & 0.440 & 0.154 & 0.0468 \\
3.0 & 0.215 & 0.614 & 0.324 & 0.152 \\
4.0 & 0.107 & 0.758 & 0.515 & 0.313 \\
5.0 & 0.051 & 0.861 & 0.682 & 0.490 \\
\bottomrule
\end{tabular}
\end{table}

The table demonstrates the central mathematical effect. For a $5k_BT_e$ cutoff, only about five percent of particles are moved in this illustrative distribution, yet the cubic energy moment falls by roughly one half. The exact bremsstrahlung kernel is not $E^3$, but higher-energy electrons contribute disproportionately, so moment reduction is a useful screening metric.

\begin{figure}[H]
\centering
\includegraphics[width=0.78\textwidth]{moment_reduction.pdf}
\caption{Radiation-weighted moment ratio versus cutoff energy for three surrogate energy exponents. More strongly energy-weighted processes benefit more from removal or redistribution of the tail.}
\label{fig:moment-reduction}
\end{figure}

\section{Interpretation for electron--ion and electron--electron emission}

Electron--ion bremsstrahlung can be represented as a one-electron average over an energy-dependent cross-section kernel. Electron--electron emission is a two-electron average and becomes more important as relativistic particles become common. The illustrative moment result suggests two qualitative expectations:
\begin{enumerate}
\item both channels decrease when high-energy electrons are moved to lower energy; and
\item electron--electron emission may respond especially strongly because encounters involving two energetic electrons are depleted.
\end{enumerate}
These expectations are consistent with relativistic calculations showing suppression under tail redistribution \cite{munirov2023}. Their magnitude must be computed using the actual cross sections and angular distribution.

\section{A surrogate spectral calculation}

To demonstrate spectral weighting, the baseline source is represented by
\begin{equation}
S_0(E_\gamma)\propto E_\gamma e^{-E_\gamma/k_BT_e},
\end{equation}
which has a bremsstrahlung-like continuum shape but is not an exact emissivity. A smooth suppression factor is applied preferentially above a cutoff-dependent photon energy. Three response functionals are then compared:
\begin{enumerate}
\item integrated energy fluence;
\item a shallow-tissue weighting that emphasizes lower photon energy; and
\item a penetrating-dose weighting that emphasizes higher photon energy.
\end{enumerate}
The weightings are intentionally dimensionless surrogates. In a real calculation they are replaced by Monte Carlo transport and mass energy-absorption coefficients.

\begin{figure}[H]
\centering
\includegraphics[width=0.78\textwidth]{dose_vs_power.pdf}
\caption{Illustrative comparison of integrated source reduction and spectrally weighted response. Different receptors need not experience the same reduction even when exposed to the same modified source.}
\label{fig:dose-vs-power}
\end{figure}

\Cref{fig:dose-vs-power} shows that the response ratio depends on the spectral weighting. In this particular constructed example, high-energy suppression produces a larger decrease in the penetrating-dose functional than in the shallow-energy-weighted functional. A different spectral suppression could reverse the ordering. The robust conclusion is not the specific curve; it is that the downstream receptor changes the benefit metric.

\section{First-order biological response}

If magnetic control changes dose but not radiation quality, oxygenation, or dose-rate response, the linear-quadratic model predicts
\begin{equation}
\frac{S_c}{S_0}
=\exp\left[-\alpha(D_c-D_0)-\beta(D_c^2-D_0^2)\right].
\label{eq:survival-ratio}
\end{equation}
Because $D_c<D_0$, the controlled-source surviving fraction is higher for normal tissue or unintended exposure. The response is nonlinear: a fifty-percent dose reduction can yield much more than a fifty-percent change in surviving fraction when the baseline dose lies on the steep part of the survival curve.

For occupational low-dose exposure, clonogenic survival is not the relevant risk endpoint. A protection-model estimate instead scales collective or individual risk approximately with effective dose under the adopted low-dose model. Thus, the same physical source reduction has application-specific biological interpretation.

\section{Illustrative foci kinetics}

For an acute exposure with initial DSB proxy $N_0=Y_DD$, single-exponential repair gives
\begin{equation}
N(t)=Y_DD e^{-k_rt}+N_{\mathrm{res}},
\label{eq:foci}
\end{equation}
where $N_{\mathrm{res}}$ represents persistent lesions or background. If source suppression only reduces dose, all time points scale approximately with $D$ after background subtraction. If the initial foci ratio equals the dose ratio but the late foci ratio differs, the spectrum or dose-rate change may have altered lesion complexity or repair. This is why time-resolved foci are more informative than a single time point.

\section{Energy-balance test}

Suppose the plasma-volume bremsstrahlung power decreases by $\Delta P_{\mathrm{plasma}}<0$, while redirected electrons produce collector radiation $\Delta P_{\mathrm{collector}}>0$. Net source suppression requires
\begin{equation}
\Delta P_{\mathrm{external}}
=\Delta P_{\mathrm{plasma}}
+\Delta P_{\mathrm{collector}}
+\Delta P_{\mathrm{secondary}}<0.
\end{equation}
For biological dose at receptor $r$, the stronger condition is
\begin{equation}
\int K_r(E,\Omega,\mathbf{x})\Delta j_{\mathrm{external}}\,
\dd E\dd\Omega\dd^3x<0.
\end{equation}
A collector can emit less total power than the original plasma while creating a more directional source near a penetration. The dose-kernel condition therefore supersedes an integrated-power condition.

\section{Sensitivity expectations}

The model predicts strong sensitivity to:
\begin{itemize}
\item cutoff energy and transition sharpness;
\item location of redistributed particles;
\item electron temperature and anisotropy;
\item impurity $Z_{\mathrm{eff}}$;
\item the fraction of lost electrons intercepted by high-$Z$ material;
\item source directionality and receptor geometry; and
\item low-energy filtration and high-energy streaming paths.
\end{itemize}
The model predicts weaker sensitivity to parameters that affect only low-energy bulk electrons when the biological receptor is dominated by penetrating hard photons. Global sensitivity analysis should test, rather than assume, this ordering.

\section{Falsifying outcomes}

The proposed mechanism would be weakened or rejected if:
\begin{enumerate}
\item the measured tail is unchanged while the detector signal falls, suggesting geometry or calibration effects;
\item electron--electron or collector emission compensates for the electron--ion reduction;
\item synchrotron power increases by more than bremsstrahlung decreases;
\item the residual spectrum hardens enough that receptor dose does not decrease;
\item equal-dose biological response differs in a way inconsistent with measured spectrum and pulse structure; or
\item perturbation-induced confinement loss reduces fusion output more than the radiative saving.
\end{enumerate}
A useful thesis must make these failure conditions explicit.

\section{Summary}

The illustrative calculations establish three robust trends: tail redistribution can strongly reduce energy-weighted moments; the benefit is larger for more strongly energy-weighted channels; and biological dose reduction is a receptor-weighted spectral quantity rather than a synonym for total-power reduction. These results justify the full cross-section, transport, and validation program.
